\begin{abstract}
To an $n\times n$ magic square we attach the $n^2$ points $(x,y,z)$ in $\R^3$, where $(x,y)$ is the position of a cell measured from the centre of the grid and $z$ is its entry minus the mean entry; their convex hull is the \emph{Magic Gem}.
Its lower and upper faces are the regular subdivisions of the grid induced by the entries and by their negatives.
For the Siamese (de la Loub\`ere) square of every odd order $n\ge5$ we prove that the gem has exactly the twelve vertices that the boundary of the grid forces (the four corners and, on each side, the cells of the smallest and largest entry), with $f$-vector $(12,26,16)$ and isometry group $\{\pm\mathrm{id}\}$; the proof reduces to checking the signs of a table of polynomials in $m=(n-1)/2$.
For the $880$ magic squares of order four we compute all gems exactly, by computer: no gem has an isometry that its point set lacks; $48$ gems are centrally symmetric, $304$ have a single half-turn symmetry and $528$ have none; the vertex set determines the square; and the gems form $616$ congruence classes and $352$ combinatorial types.
A short covariance argument shows that, for semi-magic squares of every order, the point set has a non-trivial isometry exactly when the square equals its complement up to rotation and reflection.
Other linear magic squares can have many more vertices: up to $18$, $25$, $31$ and $36$ for orders $5$, $7$, $9$ and $11$.
\medskip

\noindent\textbf{Keywords:} magic squares, convex polytopes, regular subdivisions, symmetry groups, Siamese method

\noindent\textbf{MSC 2020:} 05B15, 52B05, 52B15, 52B20
\end{abstract}
